% 原书第 15—21 页；导言、第 2.1 节。
\chapter{因果推断的假设}
\label{ch:assumptions-causal-inference}

我们已经接触了\textbf{结构因果模型（structural causal model，SCM）}的基本组成部分。现在不妨稍作停顿，考察前文涉及的一些假设，并思考这些假设对因果推理和因果学习意味着什么。

本章讨论的一个关键概念是某种形式的独立性。我们可以借助一种称为 Beuchet 椅的视错觉，先对它作直观说明。当我们看到图 2.1 左图那样的物体时，大脑会假设：物体本身与其反射光中的信息传入大脑所依赖的机制相互独立。然而，从某个极为特定的视点观察，就会违反这一假设。此时，知觉会发生错误：我们会看到一把具有三维结构的椅子，而现实中并不存在这样一把椅子。不过，在大多数情况下，独立性假设确实成立。观察物体时，大脑假定物体独立于观察位置和照明条件。因此，我们不应遇到概率极低的巧合：彼此分离的三维结构不应恰好在二维投影中对齐，阴影边界也不应恰好与纹理边界重合。在视觉研究中，这称为\textbf{一般视点假设（generic viewpoint assumption）}［Freeman, 1994］。

但独立性假设的适用范围远不止于此。下面第 2.1 节将说明，因果生成过程由若干自主模块组成，这些模块既不向彼此提供信息，也不相互影响。具体来说，虽然一个模块的输出可能影响另一个模块的输入，但模块本身彼此独立。

在上述例子中，整体知觉是物体、照明和视点的函数；然而，我们改变观察位置时，物体与照明本身并不受影响。换言之，整个因果生成模型中的某些组成部分保持不变，而我们可以利用这种不变性推断三维信息。这就是\textbf{运动恢复结构（structure from motion）}的基本思想［Ullman, 1979］，它在生物视觉和计算机视觉中都占据核心地位。

\begin{figure}[htbp]
\centering
\begin{tikzpicture}[x=1cm,y=1cm,line join=round,line cap=round]
% 左：普通视点下可见前后分离的部件；右：投影恰好对齐。
\begin{scope}
\fill[gray!10] (-.1,-.2)--(5.5,-.2)--(4.4,3.4)--(.8,3.4)--cycle;
\draw[gray!35] (-.1,-.2)--(5.5,-.2);
\begin{scope}[yshift=.4cm]
\filldraw[fill=gray!30,draw=gray!60] (1.05,2.2)--(2.8,2.35)--(3.3,1.93)--(1.45,1.75)--cycle;
\draw[line width=2.2pt,draw=brown!65!black] (2.25,2.29)--(2.32,4.18)--(2.95,4.07)--(2.87,2.06);
\draw[line width=1.8pt,draw=brown!65!black] (2.30,3.74)--(2.93,3.62);
\end{scope}
\filldraw[fill=gray!20,draw=gray!50] (.75,.1)--(3.85,.2)--(4.75,-.13)--(1.45,-.35)--cycle;
\foreach \x/\y in {1.6/.1,2.35/-.02,3.1/.18,3.7/.05}{
 \draw[line width=2.5pt,draw=brown!65!black] (\x,\y)--++(0,1.72);}
\draw[line width=2pt,draw=brown!65!black] (1.6,1.03)--(2.35,.91)--(3.7,.98)--(3.1,1.10)--cycle;
\node[font=\small,align=center] at (2.6,-.85){一般视点：各部件彼此分离};
\end{scope}
\begin{scope}[xshift=7cm]
\fill[gray!10] (-.1,-.2)--(5.2,-.2)--(4.3,3.4)--(.6,3.4)--cycle;
\draw[gray!35] (-.1,-.2)--(5.2,-.2);
\filldraw[fill=gray!20,draw=gray!50] (1.15,-.1)--(3.6,.08)--(4.12,-.35)--(1.65,-.55)--cycle;
\foreach \x/\y in {1.55/-.12,2.05/-.35,3.15/.00,3.65/-.23}{
 \draw[line width=2.5pt,draw=brown!65!black] (\x,\y)--++(0,2.12);}
\draw[line width=2pt,draw=brown!65!black] (1.55,1.07)--(2.05,.84)--(3.65,.96)--(3.15,1.19)--cycle;
\filldraw[fill=gray!30,draw=gray!60] (1.49,2.03)--(3.20,2.16)--(3.74,1.92)--(2.01,1.78)--cycle;
\draw[line width=2.2pt,draw=brown!65!black] (2.85,2.13)--(2.92,4.15)--(3.43,4.03)--(3.35,1.99);
\draw[line width=1.8pt,draw=brown!65!black] (2.90,3.69)--(3.41,3.58);
\node[font=\small,align=center] at (2.5,-.85){特殊视点：投影呈现完整椅子};
\end{scope}
\end{tikzpicture}
\caption{左图展示了从一般视点看到的 Beuchet 椅的各个（彼此分离的）部件。右图展示了从一个极为特殊的视点观察这些部件时产生的椅子错觉。从这一偶然视点看去，我们会知觉到一把椅子。原照片由 Markus Elsholz 提供；此处以 TikZ 示意重绘两种视点下的结构关系。}
\label{fig:ch02-beuchet-chair}
\end{figure}

\FloatBarrier
\section{机制独立原理}
\label{sec:independent-mechanisms}

下面先描述一个简单的原因—结果问题，并指出其中的若干现象。随后，我们将尝试统一理解这些现象之间的联系，说明它们如何源于同一个独立性原理。

假设，我们已经根据某国一组城市的样本，估计出海拔 $A$ 与年平均气温 $T$ 的联合密度 $p(a,t)$（见原书第 65 页图 4.6）。考虑以下两种表示 $p(a,t)$ 的方式：
\begin{equation}
\begin{aligned}
 p(a,t)&=p(a\mid t)\,p(t)\\
       &=p(t\mid a)\,p(a).
\end{aligned}
\tag{2.1}
\end{equation}
第一种分解描述了 $T$ 及条件分布 $A\mid T$，对应于按照图 $T\to A$ 对 $p(a,t)$ 进行因子分解。\footnote{注意，给定 $p(t)$ 后，条件密度 $p(a\mid t)$ 使我们能够计算 $p(a,t)$，进而计算 $p(a)$。暂时可以据此理解 $T\to A$ 中的箭头方向。定义 6.21 将给出严格说明。}
第二种分解则对应于按照 $A\to T$ 进行因子分解（参见定义 6.21）。我们能否判断，这两种结构中的哪一种才是因果结构？也就是说，在什么情况下，我们能够将箭头解释为因果关系？

第一个思路是考察干预的作用（见图 2.2 左侧）。设想，我们能够借助某种假想机制，抬升一座城市所处的地面，从而改变其海拔 $A$。假定我们发现，年平均气温随之下降。接下来再设想另一项干预实验：这次不改变海拔，而是在城市周围建造一个庞大的供暖系统，将年平均气温提高几度。假定我们发现，城市海拔不受影响。干预 $A$ 改变了 $T$，干预 $T$ 却没有改变 $A$。于是，我们有充分理由更倾向于用 $A\to T$ 描述因果结构。

尽管这类假想干预在实践中很难实施，甚至根本无法实施，为什么我们仍觉得上述干预效应的描述是合理的？如果改变海拔 $A$，我们假设，产生年平均气温的物理机制 $p(t\mid a)$ 仍然有效，并使 $T$ 发生变化。这里的机制包括大气的化学组成、气压如何随海拔升高而降低的物理规律，以及风的气象机制等。无论我们从哪个分布中抽取城市，这些机制都应成立，因而与 $p(a)$ 无关。奥地利人选择的建城地点可能与瑞士人略有不同，但机制 $p(t\mid a)$ 在两种情况下都适用。\footnote{这里讨论的是一种理想化情形。对于这些概括性的说法，无疑能够构造出反例。}

反过来，如果改变 $T$，就很难把 $p(a\mid t)$ 看作一个仍然有效的机制——我们大概从一开始就不相信存在这样的机制。给定一组不同的城市分布 $p(a,t)$，尽管每个分布都可以写成 $p(a\mid t)\,p(t)$，我们却会发现，无法用同一个保持不变的 $p(a\mid t)$ 解释所有分布。

这种直觉可以改述为以下两项假定：如果 $A\to T$ 是正确的因果结构，那么，
\begin{enumerate}[label=(\roman*)]
 \item 原则上可以对 $A$ 实施局部干预，即改变 $p(a)$ 而不改变 $p(t\mid a)$；
 \item $p(a)$ 与 $p(t\mid a)$ 是现实世界中具有\textbf{自主性（autonomy）}、\textbf{模块性（modularity）}或\textbf{不变性（invariance）}的机制或对象。
\end{enumerate}

有趣的是，我们虽然从假想干预实验出发来判断因果结构，推理的结果却表明，实际实施干预可能并非获知因果结构的唯一途径。对于不同的数据来源 $p(a,t)$，我们还可能通过检验式 (2.1) 的哪一种分解能产生自主或不变的组成项，识别因果结构。继续沿用上述例子，用 $p^{\ddot{\mathrm{o}}}(a,t)$ 与 $p^{\mathrm{s}}(a,t)$ 分别表示奥地利和瑞士的海拔与气温联合分布。由于两国在不同地点建城，即 $p^{\ddot{\mathrm{o}}}(a)$ 与 $p^{\mathrm{s}}(a)$ 不同，这两个联合分布也可能不同。然而，它们的因果因子分解仍可能使用相同的条件分布，即
\[
 p^{\ddot{\mathrm{o}}}(a,t)=p(t\mid a)\,p^{\ddot{\mathrm{o}}}(a),\qquad
 p^{\mathrm{s}}(a,t)=p(t\mid a)\,p^{\mathrm{s}}(a).
\]

接下来讨论另一个思路（见图 2.2 中间）。它与前例密切相关，但也有区别：它同样适用于单个分布。在因果因子分解 $p(a,t)=p(t\mid a)\,p(a)$ 中，我们预期，将条件密度 $p(t\mid a)$ 视为 $t$ 和 $a$ 的函数时，它不提供有关边缘密度函数 $p(a)$ 的信息。如果 $p(t\mid a)$ 描述的是一种不受输入分布 $p(a)$ 左右的物理机制，这一点便成立。换言之，该机制不受我们将其应用于哪一组城市的影响。

另一方面，如果写成 $p(a,t)=p(a\mid t)\,p(t)$，上述\textbf{原因与机制的独立性（independence of cause and mechanism）}便不再成立。相反，我们会发现，为了将观测到的 $p(t)$ 与 $p(a,t)$ 联系起来，机制 $p(a\mid t)$ 必须呈现出某种相当特殊的形式，并受等式 $p(a,t)=p(a\mid t)\,p(t)$ 的约束。给定一组城市及其气温数据，就可以对此进行经验检验。\footnote{第 4.1.7 节将把这一思路形式化。}

我们已经看到，独立性、自主性和不变性涉及若干不同的思路，它们都可以为因果推断提供依据。最后再考察一种思路（见图 2.2 右侧），它涉及噪声项的独立性。为便于说明，可将式 (2.1) 改写为由图为 $A\to T$ 的结构因果模型所导出的分布，并把结果 $T$ 表示为原因 $A$ 的含噪声函数：
\[
 \begin{aligned}
 A&\coloneqq N_A,\\
 T&\coloneqq f_T(A,N_T),
 \end{aligned}
\]
其中，$N_T$ 和 $N_A$ 是统计独立的噪声，即 $N_T\indep N_A$。对 $f_T$ 的函数形式施加适当限制（见第 4.1.3—4.1.6 节与第 7.1.2 节），就能够识别出，是哪一种因果结构——$A\to T$ 还是 $T\to A$——导出了观测到的 $p(a,t)$。不过，如果没有这些限制，我们始终可以实现式 (2.1) 中的两种分解。此外，在多元情形下，只要条件适当，噪声联合独立的假设便允许我们通过条件独立性检验来识别因果结构（见第 7.1.1 节）。

\begin{figure}[htbp]
\centering
\begin{tikzpicture}[causal,x=1cm,y=1cm,
 ch02box/.style={draw,rounded corners=3pt,align=center,inner sep=7pt,font=\small}]
\node[ch02box,text width=7cm] (ch02root) at (0,0)
 {\textbf{机制在物理上的独立性}\\原理 2.1};
\node[ch02box,text width=3cm] (ch02left) at (-4.6,-2.7)
 {可干预性、\\自主性、\\模块性、\\不变性、\\迁移};
\node[ch02box,text width=3cm] (ch02middle) at (0,-2.7)
 {各机制所含\\信息的独立性};
\node[ch02box,text width=3cm] (ch02right) at (4.6,-2.7)
 {噪声的独立性、\\结构的\\条件独立性};
\draw[ed] (ch02root.south west)--(ch02left.north east);
\draw[ed] (ch02root.south)--(ch02middle.north);
\draw[ed] (ch02root.south east)--(ch02right.north west);
\end{tikzpicture}
\caption{机制独立原理及其对因果推断的启示（原理 2.1）。}
\label{fig:ch02-icm}
\end{figure}


我们愿意把上述现象看作同一个一般原理的不同体现：各个机制在物理上彼此独立。

\Needspace{8\baselineskip}
\begin{principle}[独立机制]
一个系统中各变量的因果生成过程，由若干自主模块组成；这些模块既不向彼此提供信息，也不相互影响。

在概率情形下，这意味着，给定每个变量的原因时，该变量的条件分布（即其生成机制）既不提供有关其他条件分布的信息，也不影响其他条件分布。如果只有两个变量，这一原理就归结为：原因的分布与产生结果分布的机制相互独立。
\end{principle}

如果我们把系统理解为由若干模块组成，每个模块包含一个变量或一组变量，并且这些模块代表现实世界中在物理上相互独立的机制，那么这一原理就很合理。两个变量的特例被称为\textbf{原因与机制的独立性（ICM）}［Daniušis et al., 2010；Shajarisales et al., 2015］。具体来说，可以把输入视为某个制备过程的结果，而负责制备的机制独立于将输入转换为输出的机制。

在深入讨论这一原理之前，应先指出，并非所有系统都满足它。例如，如果系统的各个机制通过人为设计或演化而彼此适配，这种独立性就可能被破坏。

下面将说明，这一原理具有足够的广泛性，可以涵盖因果推理和因果学习的主要方面（见图 2.2）。我们依次讨论三个方面，分别对应图中从左到右的三条分支。

\begin{enumerate}
\item 一种理解这些模块的方式，是把它们视为具有输入—输出行为的物理机器。这一假设意味着，我们可以改变某个机制，而不影响其他机制；用因果术语来说，就是能够对一个机制实施干预，而不影响其他机制。改变某个机制会改变它的输入—输出行为，因而也会改变下游其他机制可能接收的输入；但我们假定，其他物理机制本身并不受这一改变的影响。人们往往隐含地采用这类假设，以说明干预为何可能。不过，也可以把它看作因果推理和因果学习更一般的基础。如果系统允许这种局部干预，那么各个机制之间就不存在某种由“元机制”构成、以有向方式将它们连接起来的物理通路。这也使我们有理由预期，面对所考察系统内部的变化，甚至面对某些来自系统外部的变化，各机制倾向于保持不变（见第 7.1.6 节）。这种机制的自主性预计有助于知识迁移：当一个相关领域与源领域共享部分模块时，可以把在源领域学到的知识迁移过去（见第 5.2 节和第 8.3 节）。

\item 对第一个方面的讨论侧重于独立性的物理含义及其推论，但上述论述还包含信息论方面的含义。若干彼此耦合的对象和机制在时间中的演化，能够产生统计相依关系。这与原书第 10 页的讨论相关：当时我们考察了手写数字的类别标签与图像之间的相依关系。类似地，物理上耦合的机制往往会产生信息，而这些信息可以用统计信息度量或算法信息度量来量化（见后文第 4.1.9 节和第 6.10 节）。

这里必须区分两个层次的信息。显然，结果包含关于原因的信息；但是，按照机制独立原理，将原因转换为结果的机制，不包含关于生成原因的机制的信息。对于包含两个以上节点的因果结构，该原理规定：根据各节点的直接原因生成该节点的那些机制，彼此不包含对方的信息。\footnote{这一独立性与第 1 点所述的独立性存在直观联系：如果各机制独立变化，那么一个机制的变化，就不提供有关其他机制如何变化的信息。尽管如此，第二种独立性仍包含第一种独立性并不严格涵盖的内容，因为它也适用于所有机制均未发生变化的情形；例如，它同样适用于同质数据集。}

\Needspace{5\baselineskip}
\item 最后，需要讨论结构方程建模中常用的噪声项独立假设，与机制独立原理之间有何联系。这种联系不那么直观。为此，考虑变量 $E\coloneqq f(C,N)$，其中噪声 $N$ 是离散的。对于 $N$ 的每一个取值 $s$，赋值 $E\coloneqq f(C,N)$ 就化为确定性机制 $E\coloneqq f^s(C)$，将输入 $C$ 转换为输出 $E$。实际上，这意味着噪声在若干机制 $f^s$ 中随机选择一个，机制的数量等于噪声变量 $N$ 值域的基数。现在假设，节点 $X_j$ 和 $X_k$ 对应的两个机制，其噪声变量在统计上相依。\footnote{尽管到目前为止我们主要关注两个变量的情形，这里的论证表述同样适用于包含两个以上变量的因果结构。}
例如，这种相依关系可能保证：一旦节点 $j$ 上启用了某个机制 $f_j^s$，我们便知道节点 $k$ 上启用的是哪个机制 $f_k^t$。这将违反机制独立原理。

上一段把噪声变量视为不同机制之间的选择器，是一种较为极端的理解方式（亦见第 3.4 节）。在实际问题中，噪声的作用可能并没有那么显著。例如，如果噪声是加性的，即 $E\coloneqq f(C)+N$，它对机制的影响就受到限制：它只能使机制的输出向上或向下平移，因此只是在一组彼此非常相似的机制之间进行选择。这与另一种观点一致：噪声变量是我们试图描述的系统之外的变量，反映了系统永远不可能与环境完全隔离这一事实。按照这种理解，噪声之间可能存在较弱的相依关系，而机制独立原理仍不失效。
\end{enumerate}

原理 2.1 的上述各个方面，都可能有助于因果学习；换言之，它们都可能提供有关因果结构的信息。不过，可以设想，这些信息在某些情况下也可能相互冲突，具体取决于在给定情境中哪些假设成立。
\FloatBarrier
